\documentclass[11pt]{article}
\usepackage[utf8]{inputenc}
\usepackage[T1]{fontenc}
\usepackage[ngerman]{babel}
\usepackage[margin=2.5cm]{geometry}
\usepackage{tikz}
\usetikzlibrary{arrows.meta, decorations.pathreplacing}

\title{Algorithm Engineering\\Woche 2, Aufgabe 1}
\author{Jan Vogt}
\date{}

\colorlet{tnull}{blue!35}
\colorlet{teins}{orange!55}

\def\w{0.72}
\def\h{0.6}

\newcommand{\zeile}[3]{
  \foreach \i in {0,...,7} {
    \pgfmathtruncatemacro{\idx}{#3 + \i}
    \draw[fill=white] (#1 + \i * \w, #2) rectangle ++(\w, \h);
    \node[font=\scriptsize] at (#1 + \i * \w + \w / 2, #2 + \h / 2) {a[\idx]};
  }
}

\newcommand{\feld}[5]{
  \draw[fill=#5] (#1 + #3 * \w, #2) rectangle ++(\w, \h);
  \node[font=\scriptsize\bfseries] at (#1 + #3 * \w + \w / 2, #2 + \h / 2) {a[#4]};
}

\tikzset{
  kern/.style={draw, rounded corners, minimum width=5.76cm, minimum height=0.8cm, fill=gray!10},
  nr/.style={circle, draw, fill=white, inner sep=1pt, font=\scriptsize\bfseries},
  rahmen/.style={draw, rounded corners=8pt, gray!60, line width=0.8pt},
}

\begin{document}

\maketitle

\begin{center}
\begin{tikzpicture}

  \draw[rahmen] (-0.4, 0.4) rectangle (14.56, 9.2);
  \node[font=\bfseries] at (7.08, 8.8) {Problem: False Sharing (\texttt{++a[id]})};

  \node[kern] at (2.88, 7.8) {Kern 0: Thread 0 macht \texttt{++a[0]}};
  \node[kern] at (11.28, 7.8) {Kern 1: Thread 1 macht \texttt{++a[1]}};

  \zeile{0}{5.6}{0}
  \feld{0}{5.6}{0}{0}{tnull}
  \zeile{8.4}{5.6}{0}
  \feld{8.4}{5.6}{1}{1}{teins}
  \node[font=\small, anchor=north west] at (0, 5.45) {L1-Cache von Kern 0};
  \node[font=\small, anchor=north west] at (8.4, 5.45) {L1-Cache von Kern 1};

  \draw[red, line width=1.5pt, dashed] (8.4 - 0.08, 5.6 - 0.08) rectangle ++(8 * \w + 0.16, \h + 0.16);

  \draw[-{Stealth}, thick, blue!70!black] (0.36, 7.4) -- (0.36, 6.2);
  \node[nr] at (0.75, 6.8) {1};
  \node[font=\small, anchor=west] at (1.0, 6.8) {schreibt a[0]};

  \draw[-{Stealth}, thick, red] (5.76, 6.0) to[bend left=30] node[nr, midway] {2} (8.3, 6.0);
  \node[font=\small, red, align=center] at (7.08, 7.05) {macht Kopie\\ ungültig};

  \draw[rounded corners, fill=gray!10] (3.4, 1.0) rectangle (10.76, 3.0);
  \node[font=\small] at (7.08, 2.7) {gemeinsamer Speicher (L3-Cache / RAM)};
  \zeile{4.2}{1.6}{0}
  \feld{4.2}{1.6}{0}{0}{tnull}
  \feld{4.2}{1.6}{1}{1}{teins}
  \draw[decorate, decoration={brace, mirror, amplitude=4pt}]
    (4.2, 1.5) -- (4.2 + 8 * \w, 1.5) node[midway, below=4pt, font=\scriptsize] {eine Cache Line = 64 Byte};

  \draw[-{Stealth}, thick, red] (7.9, 3.0) -- (8.25, 5.45) node[nr, pos=0.45] {3};
  \node[font=\small, red, align=left, anchor=west] at (8.5, 4.2) {Kern 1 muss die ganze\\ Cache Line neu laden};

  \draw[-{Stealth}, thick, gray] (5.0, 3.0) -- (4.6, 5.45);

  \node[font=\small, align=center] at (1.6, 3.9) {danach dasselbe\\ andersrum, die Line\\ wandert ständig\\ hin und her};

  \draw[rahmen] (-0.4, -7.4) rectangle (14.56, -0.2);
  \node[font=\bfseries] at (7.08, -0.8) {Lösung: Padding (\texttt{++a[id * 8]})};

  \node[kern] at (2.88, -1.8) {Kern 0: Thread 0 macht \texttt{++a[0]}};
  \node[kern] at (11.28, -1.8) {Kern 1: Thread 1 macht \texttt{++a[8]}};

  \zeile{0}{-4.0}{0}
  \feld{0}{-4.0}{0}{0}{tnull}
  \zeile{8.4}{-4.0}{8}
  \feld{8.4}{-4.0}{0}{8}{teins}
  \node[font=\small, anchor=north west] at (0, -4.05) {L1-Cache von Kern 0};
  \node[font=\small, anchor=north west] at (8.4, -4.05) {L1-Cache von Kern 1};

  \draw[-{Stealth}, thick, blue!70!black] (0.36, -2.2) -- (0.36, -3.4);
  \node[font=\small, anchor=west] at (0.6, -2.8) {schreibt a[0]};
  \draw[-{Stealth}, thick, orange!80!black] (8.76, -2.2) -- (8.76, -3.4);
  \node[font=\small, anchor=west] at (9.0, -2.8) {schreibt a[8]};

  \node[font=\small, green!50!black, align=center] at (7.08, -3.0) {kein\\ Konflikt};

  \draw[rounded corners, fill=gray!10] (0.8, -7.1) rectangle (13.36, -4.85);
  \zeile{1.32}{-6.1}{0}
  \feld{1.32}{-6.1}{0}{0}{tnull}
  \zeile{7.08}{-6.1}{8}
  \feld{7.08}{-6.1}{0}{8}{teins}
  \draw[decorate, decoration={brace, mirror, amplitude=4pt}]
    (1.32, -6.2) -- (1.32 + 8 * \w, -6.2) node[midway, below=4pt, font=\scriptsize] {Cache Line 1};
  \draw[decorate, decoration={brace, mirror, amplitude=4pt}]
    (7.08, -6.2) -- (7.08 + 8 * \w, -6.2) node[midway, below=4pt, font=\scriptsize] {Cache Line 2};
  \node[font=\small] at (7.08, -5.2) {gemeinsamer Speicher (L3-Cache / RAM)};

\end{tikzpicture}
\end{center}

\subsection*{Messung}

Mit \texttt{false\_sharing\_test.cpp} habe ich das Programm von Folie 9 und 10 mit 4 Threads und
verschiedenen Abständen zwischen den Zählern laufen lassen. Jede Variante lief 10 mal, weil die
Zeiten auf meinem Laptop (Apple M3) ziemlich stark schwanken.

\begin{center}
\begin{tabular}{rrrr}
PAD & Abstand & beste Zeit [s] & Median [s] \\
\hline
1 & 8 Byte & 1.200 & 1.408 \\
8 & 64 Byte & 0.807 & 1.040 \\
16 & 128 Byte & 0.800 & 1.127 \\
\end{tabular}
\end{center}

Ohne Padding ist das Programm also deutlich langsamer, bei der besten Zeit ungefähr um den Faktor
1,5. Laut \texttt{sysctl hw.cachelinesize} hat der M3 eigentlich 128 Byte große Cache Lines.
Trotzdem war bei meiner Messung zwischen 64 und 128 Byte Abstand kein klarer Unterschied zu
sehen. Woran das genau liegt, konnte ich nicht herausfinden.

\end{document}
